\documentclass{article}
\usepackage{PRIMEarxiv} % Core package for the PrimeAI Template
\usepackage{amsmath, amssymb, amsthm, bm, dcolumn} % Add amsthm here for the proof environment
\usepackage[numbers,sort&compress]{natbib} % Natbib for citations
\usepackage{graphicx} % For high-quality images
\usepackage{multicol} % Multi-column figures/tables
\usepackage{hyperref} % Hyperlinks
\usepackage{listings} % Code listings
\usepackage{authblk} % For structured affiliations
% Define theorem style (optional)
\theoremstyle{plain}
\newtheorem{theorem}{Theorem}
\renewcommand{\qedsymbol}{}

% Custom commands for primes and qubit encoding
\newcommand{\primeQubit}[1]{|p_{#1}\rangle}
\newcommand{\primeGate}[1]{U_{p_{#1}}}

\usepackage{wrapfig}
\usepackage[pscoord]{eso-pic}
\usepackage[fulladjust]{marginnote}
\reversemarginpar

% Typesetting improvements without footnote patching
\usepackage[protrusion=true, expansion=true, tracking=false]{microtype}
\microtypecontext{spacing=nonfrench}

% Line numbers
\usepackage[right]{lineno}

% Text layout - adjust as needed
\raggedright
\setlength{\parindent}{0.5cm}
\textwidth 5.25in 
\textheight 8.75in

% Set double spacing
\usepackage{setspace} 
\doublespacing

% Adjust width for specific content
\usepackage{changepage}

% Adjust caption style
\usepackage[aboveskip=1pt,labelfont=bf,labelsep=period,singlelinecheck=off]{caption}

% Remove brackets from references
\makeatletter
\renewcommand{\@biblabel}[1]{\quad#1.}
\makeatother

% Header, footer, and page numbers
\usepackage{lastpage,fancyhdr}
\pagestyle{fancy}
\fancyhf{}
\fancyhead[L]{Preprint - PrimeAI Enhanced Template}
\fancyfoot[C]{\scriptsize Multiplicity Theory © 2024 Ryan Van Gelder - Citizen Gardens \\ Licensed Under MIT and CC BY-NC-SA 4.0.}
\fancyfoot[R]{Page \thepage\ of \pageref{LastPage}}
\renewcommand{\footrule}{\hrule height 2pt \vspace{2mm}}

\section{Secure Quantum Communication and Cryptography}

Advancements in secure quantum communication rely on entanglement-based encryption, prime-encoded error correction, and resilient cryptographic methodologies. By leveraging Fibonacci/Lucas sequences and prime-indexed encoding, these quantum-AI-driven security frameworks ensure robust, fault-tolerant cryptographic mechanisms.

\subsection{Entanglement-Based Quantum Cryptographic Security}

Quantum key exchange (QKE) using entanglement enhances secure communication channels by ensuring that encryption keys remain synchronized across long distances. The entanglement-based cryptographic state is defined as:

\begin{equation}
|\Psi\rangle = \frac{1}{\sqrt{2}} (|00\rangle + |11\rangle),
\end{equation}

where:
\begin{itemize}
    \item $|\Psi\rangle$ is the Bell state entangling two remote quantum nodes.
    \item Secure key exchange is achieved when measurements on both nodes produce correlated results.
    \item Any interception collapses the entangled state, alerting communicating parties.
\end{itemize}

Applications:
\begin{itemize}
    \item Quantum-secure authentication protocols for cryptographic networks.
    \item Non-clonable key exchange mechanisms resistant to computational attacks.
    \item Entanglement-enhanced blockchain security for decentralized cryptographic ledgers.
\end{itemize}

\subsection{Quantum Error Correction with Prime-Encoding}

Quantum error correction (QEC) using Fibonacci/Lucas-based prime encoding enhances resilience against decoherence and computational noise. The error-corrected quantum state follows:

\begin{equation}
E_{\text{corrected}} = E_{\text{measured}}(1 - \Phi^{-2}),
\end{equation}

where:
\begin{itemize}
    \item $E_{\text{measured}}$ is the detected error probability.
    \item $\Phi = \frac{1+\sqrt{5}}{2}$ represents the golden ratio-based error suppression factor.
    \item $\Phi^{-2}$ ensures non-linear recursive stability in quantum state correction.
\end{itemize}

Advantages:
\begin{itemize}
    \item Fibonacci/Lucas-weighted QEC reduces error propagation in quantum computations.
    \item Self-correcting AI-driven quantum networks ensure real-time resilience.
    \item Prime-weighted noise mitigation minimizes entropy fluctuations in quantum encryption.
\end{itemize}

\subsection{Quantum Resilience in Error Correction}

Entropy stabilization in quantum cryptographic environments requires prime-indexed reinforcement mechanisms. The entropy-corrected state follows:

\begin{equation}
S_{\text{prime}} = \sum_{i} p_i \log_2 p_i,
\end{equation}

where:
\begin{itemize}
    \item $S_{\text{prime}}$ represents the entropy function regulated by prime indices.
    \item $p_i$ denotes prime-weighted probability distributions over quantum states.
    \item Logarithmic scaling maintains stability against probabilistic fluctuations.
\end{itemize}

Implementation benefits:
\begin{itemize}
    \item Ensures stable quantum key distribution against adversarial noise.
    \item Regulates entropic complexity in AI-driven cryptographic frameworks.
    \item Reduces uncertainty in prime-based cryptographic hashing functions.
\end{itemize}
\subsection{Prime-Encoded Quantum Cryptographic Key Generation}

Secure cryptographic key generation is achieved through prime-weighted Fibonacci sequences. The key $K_n$ is derived as:
\begin{equation}
K_n = \text{SHA256} \left( \prod_{i=1}^{N} p_i^{F_i} \right),
\end{equation}
where $p_i$ represents the $i$-th prime, and $F_i$ is the $i$-th Fibonacci number. This approach ensures:
\begin{itemize}
    \item Non-reversible key structures due to prime-based multiplicative complexity.
    \item Enhanced entropy distribution across cryptographic keys.
    \item Resistance to classical and quantum brute-force attacks via prime-weighted key derivations.
\end{itemize}
\subsection{Prime-Twisted Quantum Encryption}

Prime-twisted encryption applies recursive prime-weighted transformations for secure cryptographic encoding. The encryption function follows:

\begin{equation}
K_{\text{secure}}(t) = H \left(\prod_{i=1}^{N} p_i^{F_i} \right),
\end{equation}

where:
\begin{itemize}
    \item $K_{\text{secure}}(t)$ is the dynamically generated secure cryptographic key.
    \item $H(\cdot)$ represents a secure cryptographic hash function.
    \item $p_i^{F_i}$ encodes prime-weighted Fibonacci transformations.
\end{itemize}

Security features:
\begin{itemize}
    \item Resistant to quantum brute-force decryption.
    \item Fault-tolerant key generation using prime-modulated recursive encoding.
    \item Enables adaptive, self-evolving cryptographic security layers.
\end{itemize}

\nocite{*}
\bibliographystyle{unsrt}
\bibliography{references}

\end{document}